% =====================================================================
%  Chapter 1: Digital educational systems and learning management
%             systems (LMS)
% =====================================================================
\chapter{Digital Educational Systems and Learning Management Systems}
\label{chap-lms}

\section*{Introduction}
\addcontentsline{toc}{section}{Introduction}

Higher education now relies on digital environments that organize the
delivery of courses, the exchanges between participants and the monitoring
of learners. At the centre of these environments are learning management
systems (LMSs), which have become the basic infrastructure of most
institutions \cite{dahlstrom2014,coates2005}. Our work aims to extend such
a system with a multi-agent framework able to adapt learning to each
learner and to their context. Before presenting this contribution, we need
to understand what LMSs do, how they are built and why, on their own, they
cannot personalize learning.

This chapter first defines online learning and its modalities
(\cref{sec-elearning}). It then presents LMSs, their functions and their
general architecture (\cref{sec-lms}), followed by the standards that let
them communicate with other tools (\cref{sec-standards}). It compares the
main platforms on the market (\cref{sec-overview}) and studies Moodle in
detail, since it is the platform chosen to integrate our framework
(\cref{sec-moodle}). It ends with an analysis of the limits of current
LMSs with respect to adaptive and personalized learning
(\cref{sec-limits}).

% =====================================================================
\section{Online learning}
\label{sec-elearning}
% =====================================================================

\subsection{Definitions}

The term \foreign{e-learning} covers many realities, and the literature
often uses it interchangeably with distance education or online learning.
\citet{moore2011} show that these terms are not equivalent: distance
education refers first of all to the physical separation between teacher
and learner, whereas online learning refers to the use of the Internet to
access content and interactions, whether or not the learner is remote.

\begin{definitionbox}{Online learning}
We call \emph{online learning} (or \foreign{e-learning}) any learning
situation in which content, activities, interactions or assessment are
mediated by a connected digital system, in particular a web platform
\cite{clark2016,moore2011}.
\end{definitionbox}

The broader term \emph{technology-enhanced learning environment} refers to
any computer system designed to support human learning. An LMS is a
particular case of such an environment: it focuses on organizing and
managing teaching rather than on fine-grained pedagogical interaction with
each learner.

\subsection{Evolution of distance education}

Distance education has gone through several generations, each tied to a
dominant communication technology. \citet{taylor2001} identifies five, from
correspondence teaching to automated, interactive online environments.
\Cref{fig-generations} extends this view up to current systems that
integrate generative artificial intelligence.

\begin{figure}[htbp]
\centering
\adjustbox{max width=\textwidth}{%
\begin{tikzpicture}[
  gen/.style={bloc,text width=2.35cm,minimum height=2.6cm,font=\footnotesize},
  genai/.style={gen,fill=thprimary,text=white},
  node distance=0.35cm]
  \node[gen]               (g1) {\textbf{Corres\-pondence}\\[3pt]print,\\postal mail};
  \node[gen,right=of g1]   (g2) {\textbf{Multimedia}\\[3pt]radio, television,\\cassettes};
  \node[gen,right=of g2]   (g3) {\textbf{Tele-learning}\\[3pt]video and audio\\conferencing};
  \node[gen,right=of g3]   (g4) {\textbf{Online learning}\\[3pt]web, LMS,\\forums};
  \node[gen,right=of g4]   (g5) {\textbf{Open platforms}\\[3pt]MOOCs, mobile,\\analytics};
  \node[genai,right=of g5] (g6) {\textbf{AI-augmented}\\\textbf{platforms}\\[3pt]LLMs, agents,\\adaptation};
  \draw[arrow] ([yshift=-0.35cm]g1.south west) -- ([yshift=-0.35cm]g6.south east)
     node[pos=0,below right,font=\footnotesize,text=thgray]{more interaction, finer monitoring}
     node[pos=1,below left,font=\footnotesize,text=thgray]{time};
\end{tikzpicture}}
\caption[Generations of distance education]{Generations of distance
education, after \citet{taylor2001} and Moore and Kearsley~\cite{moore2012},
extended with AI-augmented platforms.}
\label{fig-generations}
\end{figure}

\subsection{Learning modalities}

Learning arrangements differ along two axes: the share of face-to-face
teaching and the synchronicity of exchanges \cite{moore2012}.

\begin{itemize}
  \item \textbf{Web-enhanced face-to-face}: classes take place in a
        classroom, and the platform is used to share resources and to
        monitor progress.
  \item \textbf{Blended}: part of the activities takes place online and
        replaces part of the face-to-face teaching. Garrison and
        Kanuka~\cite{garrison2004} stress that this modality requires
        rethinking course design and cannot be reduced to adding digital
        resources.
  \item \textbf{Fully online}: all activities are mediated by the platform.
  \item \textbf{Synchronous or asynchronous}: exchanges happen in real
        time (virtual classroom) or at different times (forum, assignment,
        quiz).
\end{itemize}

In Algeria, the Covid-19 pandemic forced universities to move quickly to
distance teaching in 2020. They relied on several tools, including Moodle,
Google Classroom and videoconferencing services \cite{kerrouzi2023}. Since
then, blended teaching has become established in many institutions, and
Moodle has remained the most widespread institutional platform.

\subsection{Stakeholders in online learning}
\label{subsec-stakeholders}

Four categories of stakeholders interact in an online learning
arrangement (\cref{tab-stakeholders}). They will reappear in the design of
our framework (\cref{chap-contribution}).

\begin{table}[htbp]
\centering
\caption{Stakeholders in online learning and their needs}
\label{tab-stakeholders}
\begin{tabularx}{\textwidth}{@{}L{2.8cm}YY@{}}
\toprule
\textbf{Stakeholder} & \textbf{Role} & \textbf{Main needs} \\
\midrule
Learner & Follows courses, completes activities and assessments &
  A path adapted to their level, explanations, quick feedback, help
  organizing their work \\
Teacher & Designs the course, leads it, assesses learners &
  Time savings, an overview of class progress, early detection of
  difficulties, better course materials \\
Academic administrator & Organizes courses and programmes &
  Aggregated indicators, programme quality and consistency, detection of
  students at risk \\
Technical administrator & Installs and maintains the platform &
  Security, performance, integration with the information system,
  control over data \\
\bottomrule
\end{tabularx}
\end{table}

% =====================================================================
\section{Learning management systems}
\label{sec-lms}
% =====================================================================

\subsection{Definition}

\begin{definitionbox}{LMS}
A \emph{learning management system} (LMS) is an application, usually
web-based, used to administer, deliver, track and assess learning
activities. It manages users, courses, enrolments, resources, activities,
grades and communication between participants \cite{watson2007,ellis2009}.
\end{definitionbox}

Watson and Watson~\cite{watson2007} point out the frequent confusion
between LMSs and neighbouring systems. \Cref{tab-lms-neighbours} clarifies
these distinctions, which help us state precisely the role of the platform
in our architecture.

\begin{table}[htbp]
\centering
\caption{LMSs and neighbouring systems}
\label{tab-lms-neighbours}
\begin{tabularx}{\textwidth}{@{}L{3.2cm}YY@{}}
\toprule
\textbf{System} & \textbf{Main function} & \textbf{Difference from an LMS} \\
\midrule
LMS & Manage teaching: users, courses, activities, grades, tracking &
  Reference \\
LCMS (\foreign{Learning Content Management System}) & Author, store and
  reuse learning content & Focused on content rather than on organizing
  and tracking learners \\
VLE (\foreign{Virtual Learning Environment}) & Provide an online course
  space & Mostly British term, often a synonym of LMS in higher
  education \\
MOOC platform & Deliver open courses to very large numbers of learners &
  Open enrolment, little individual support, heavily automated
  assessment \\
Intelligent tutoring system (ITS) & Guide the learner step by step through
  problem solving & Detailed learner and domain models, but limited to one
  domain \\
\bottomrule
\end{tabularx}
\end{table}

\subsection{Functions}

The functions of an LMS fall into six groups \cite{ellis2009,dobre2015}:

\begin{enumerate}
  \item \textbf{User and role management}: accounts, authentication, roles
        (learner, teacher, manager) and their permissions.
  \item \textbf{Course management}: course creation, organization into
        sections, enrolments, groups, calendar.
  \item \textbf{Content delivery}: files, pages, videos, SCORM or H5P
        packages, links to external resources.
  \item \textbf{Activities and assessment}: assignments, quizzes, forums,
        peer-assessment workshops, gradebook, competencies.
  \item \textbf{Communication and collaboration}: messaging,
        announcements, forums, wikis, virtual classrooms.
  \item \textbf{Tracking and reporting}: activity logs, activity
        completion, participation reports, dashboards.
\end{enumerate}

In recent years, learning analytics modules have been added to these core
functions. Learning analytics is the collection and analysis of learners'
traces in order to understand and improve learning \cite{siemens2011}.
The first generative AI tools have also appeared
(\cref{subsec-moodle-ai}).

\subsection{General architecture}
\label{subsec-lms-architecture}

Most LMSs follow a layered web architecture, shown in
\cref{fig-lms-architecture}. A presentation layer (web browser, mobile
app) communicates with an application server that holds the business
logic. The server relies on a relational database and on a file store. An
integration layer exposes web services and implements the interoperability
standards, which connects the LMS to the institution's information system
and to external tools. Our framework will use this integration layer.

\begin{figure}[htbp]
\centering
\adjustbox{max width=\textwidth}{%
\begin{tikzpicture}[node distance=0.45cm and 0.45cm,
  b/.style={bloc,text width=2.6cm,minimum height=0.95cm}]
  % Presentation
  \node[b] (nav) {Web browser};
  \node[b,right=of nav] (mob) {Mobile app};
  \node[b,right=of mob] (ext) {External tools};
  \node[layer,fit=(nav)(mob)(ext),
        label={[font=\footnotesize\bfseries,text=thaccent]left:Presentation}] (cP) {};
  % Application
  \node[b,below=1.3cm of nav] (u) {Users, roles, permissions};
  \node[b,right=of u] (c) {Courses, activities, assessment};
  \node[b,right=of c] (s) {Tracking, logs, reports};
  \node[layer,fit=(u)(c)(s),
        label={[font=\footnotesize\bfseries,text=thaccent]left:Application}] (cA) {};
  % Integration (right)
  \path (s.east) -- (s.east |- ext.east) coordinate[pos=0.5] (mid);
  \node[blocstrong,text width=2.8cm,minimum height=4.0cm,anchor=west]
        (int) at ([xshift=1.1cm]mid) {Integration layer\\[4pt]\mdseries web services\\REST API\\LTI, xAPI,\\SCORM, QTI};
  % Data
  \node[b,below=1.3cm of u] (db) {Relational database};
  \node[b,right=of db] (fs) {File storage};
  \node[b,right=of fs] (lg) {Event logs};
  \node[layer,fit=(db)(fs)(lg),
        label={[font=\footnotesize\bfseries,text=thaccent]left:Data}] (cD) {};
  \draw[arrow,<->] (cP.south) -- (cA.north);
  \draw[arrow,<->] (cA.south) -- (cD.north);
  \draw[arrow,<->] (cA.east) -- (int.west |- cA.east);
  \draw[arrow,<->] (cP.east) -- (int.west |- cP.east);
\end{tikzpicture}}
\caption{Layered architecture of an LMS}
\label{fig-lms-architecture}
\end{figure}

\subsection{Licensing and deployment models}

From an institution's point of view, two criteria distinguish LMSs:

\begin{itemize}
  \item \textbf{Licence}: open-source LMSs (Moodle, Open edX, Chamilo,
        Sakai) can be installed, studied and modified without licence
        fees; proprietary LMSs (Blackboard Learn, D2L Brightspace) are
        provided under contract. Canvas sits in between: its core code is
        open, but it is mostly used as a service hosted by its vendor.
  \item \textbf{Deployment}: the institution can host the platform on its
        own servers, which gives it control over the data, or use it as
        Software as a Service (SaaS), which reduces the operating burden
        but places the data with a provider.
\end{itemize}

For a research project that must read learning traces and add external
services, an open-source, self-hosted platform gives full access to the
data and to the extension points.

% =====================================================================
\section{Interoperability standards}
\label{sec-standards}
% =====================================================================

An LMS never works alone. It must import content produced elsewhere,
launch external tools, exchange grades and send activity traces. Open
standards govern these exchanges. They are published mainly by the
Advanced Distributed Learning (ADL) initiative, by the 1EdTech consortium
(formerly IMS Global Learning Consortium) and by the IEEE.

\subsection{SCORM}

SCORM (\foreign{Sharable Content Object Reference Model}) defines how to
package learning content and how this content communicates with the LMS
that launches it, for example to report a score or a completion status
\cite{adl_scorm}. Its latest version, SCORM 2004 4th Edition, dates from
2009. SCORM is still very widespread, but it assumes that the content runs
in the browser, inside the LMS, and it describes only a limited amount of
information about the learner.

\subsection{xAPI and cmi5}

The Experience API (xAPI) removes these limits. It describes each learning
experience as an \enquote{actor, verb, object} statement, enriched with a
result and a context, and stored in a dedicated repository called a
Learning Record Store (LRS). A statement can come from an LMS, a mobile
app, a simulator or any other source. After several years of
standardization, xAPI became the IEEE 9274.1.1-2023 standard, sometimes
called xAPI 2.0 \cite{ieee_xapi,adl_xapi2023}.

The cmi5 specification is an xAPI profile that sets the rules for
launching content from an LMS and the vocabulary of the statements. It
thus covers what SCORM did, on top of xAPI \cite{adl_cmi5}.

The \texttt{context} field of xAPI statements is directly relevant to our
work: it carries information about the learning situation (instructor,
group, platform, parent activity) that can feed a context model.

\subsection{LTI}

The 1EdTech LTI standard (\foreign{Learning Tools Interoperability}) lets
an LMS, called the \emph{platform}, launch an external application, called
the \emph{tool}, while securely passing on the user's identity, role and
originating course \cite{lti13}. Version 1.3, published in 2019, relies on
OAuth~2.0 and OpenID Connect with signed tokens (JWT). It is complemented
by the LTI Advantage services: sending grades back to the gradebook
(\foreign{Assignment and Grade Services}), accessing the list of
participants (\foreign{Names and Role Provisioning Services}) and
selecting content (\foreign{Deep Linking}).

LTI 1.3 is the most suitable standard for our framework: it embeds an
external interface in a Moodle course without modifying the core of the
platform, while passing on the launch context.

\subsection{QTI}

The QTI standard (\foreign{Question and Test Interoperability}) describes
questions and tests in an XML format so that they can be exchanged between
platforms \cite{qti3}. It matters for our work because questions generated
by an assessment agent will have to be imported into a question bank.

\subsection{Comparative summary}

\Cref{tab-standards} compares these standards by role and by their
relevance to an external adaptive learning framework.

\begin{table}[htbp]
\centering
\caption{Comparison of LMS interoperability standards}
\label{tab-standards}
\small
\begin{tabularx}{\textwidth}{@{}L{1.9cm}L{2cm}YY@{}}
\toprule
\textbf{Standard} & \textbf{Body} & \textbf{Role} & \textbf{Relevance to our framework} \\
\midrule
SCORM 2004 & ADL & Package content and track its completion in the LMS &
  Low: little information about the learner \\
xAPI (IEEE 9274.1.1) & ADL, IEEE & Record varied learning experiences in
  an LRS & High: rich traces with a context field \\
cmi5 & ADL & Launch content from an LMS on top of xAPI &
  Medium: useful if generated content is delivered as packages \\
LTI 1.3 and Advantage & 1EdTech & Launch an external tool securely, return
  grades & High: embeds the framework's interface in Moodle \\
QTI 3.0 & 1EdTech & Exchange questions and tests &
  Medium: export of generated questions \\
\bottomrule
\end{tabularx}
\end{table}

% =====================================================================
\section{Overview of the main LMSs}
\label{sec-overview}
% =====================================================================

A small number of platforms dominate the LMS market. In North America,
Canvas, D2L Brightspace, Blackboard Learn and Moodle have held the top four
places for some fifteen years. At the end of 2024, measured by student
enrolment, Canvas held 50\,\% of the market, ahead of Brightspace
(20\,\%), Blackboard (12\,\%) and Moodle (9\,\%) \cite{onedtech2024}.
Almost all new contracts go to Canvas and Brightspace, while Moodle and
Blackboard are slowly declining in the United States
\cite{edutechnica2025}. The picture is very different worldwide: Moodle
has more than 160,000 registered sites in 242 countries and more than 400
million user accounts, and these figures only cover sites that register
voluntarily \cite{moodle_stats}. The most represented countries are Spain,
the United States, Germany, Mexico and Brazil.

\subsection{The platforms}

\paragraph{Moodle.} Open-source platform written in PHP, created by Martin
Dougiamas and released in 2002 \cite{dougiamas2003}. It is grounded in a
social constructivist approach and built around plugins. It is studied in
detail in \cref{sec-moodle}.

\paragraph{Canvas LMS.} Developed by Instructure, founded in 2008. Canvas
is valued for its modern interface and its web-service-oriented design. It
is now the most used platform in North American higher education
\cite{onedtech2024}.

\paragraph{Blackboard Learn.} The market pioneer, founded in 1997 and now
published by Anthology. The platform offers two experiences,
\enquote{Original} and \enquote{Ultra}, and has been losing market share
for several years \cite{edutechnica2025}.

\paragraph{D2L Brightspace.} Proprietary platform from the Canadian
company D2L, founded in 1999. It emphasizes accessibility,
competency-based learning and analytics.

\paragraph{Open edX.} Open-source platform from the edX project of MIT and
Harvard, opened in 2013. It is designed for massive open online courses
(MOOCs) rather than for managing a full university curriculum.

\paragraph{Google Classroom.} Online service launched by Google in 2014
and integrated into Google Workspace. It is easy to adopt and covers the
distribution of assignments and communication, but it remains limited in
advanced assessment, tracking and personalization.

\paragraph{Chamilo.} Open-source LMS launched in 2010, widely used in the
Spanish- and French-speaking worlds. It is lighter than Moodle and targets
organizations that favour simplicity.

\subsection{Comparison}

\Cref{tab-lms-comparison} compares these platforms against the criteria
that matter for our project: open code, self-hosting, extensibility,
support for standards, access to traces and built-in AI tools.

\begin{table}[htbp]
\centering
\caption[Comparison of the main LMSs]{Comparison of the main LMSs against
criteria relevant to integrating an external framework (\cyes: yes,
\cpartly: partly, \cno: no)}
\label{tab-lms-comparison}
\small
\begin{tabularx}{\textwidth}{@{}L{2.4cm}C{1.3cm}C{1.4cm}C{1.3cm}C{1.3cm}C{1.4cm}Y@{}}
\toprule
\textbf{LMS} & \textbf{Open code} & \textbf{Self-hosting} & \textbf{REST API} &
  \textbf{LTI 1.3} & \textbf{Access to traces} & \textbf{Main audience} \\
\midrule
Moodle          & \cyes    & \cyes    & \cyes    & \cyes    & \cyes    & Universities, schools, companies \\
Canvas LMS      & \cpartly & \cpartly & \cyes    & \cyes    & \cpartly & Higher education \\
Blackboard Learn& \cno     & \cpartly & \cyes    & \cyes    & \cpartly & Higher education \\
D2L Brightspace & \cno     & \cno     & \cyes    & \cyes    & \cpartly & Higher and continuing education \\
Open edX        & \cyes    & \cyes    & \cyes    & \cyes    & \cyes    & MOOCs \\
Google Classroom& \cno     & \cno     & \cyes    & \cpartly & \cno     & Schools, small groups \\
Chamilo         & \cyes    & \cyes    & \cpartly & \cpartly & \cpartly & Organizations, training \\
\bottomrule
\end{tabularx}
\end{table}

\todonote{Check the \cpartly{} cells of the table against each platform's
official documentation before submission.}

Moodle is the only platform that combines fully open code, self-hosting,
full access to logs and complete support for the standards. It is also the
platform already in place in most Algerian universities. These two reasons
justify choosing it as the integration and validation platform for our
framework.

% =====================================================================
\section{A detailed study of Moodle}
\label{sec-moodle}
% =====================================================================

\subsection{History and versions}

Moodle (\foreign{Modular Object-Oriented Dynamic Learning Environment})
grew out of Martin Dougiamas's doctoral work, which sought a platform that
fosters the collective construction of knowledge \cite{dougiamas2003}.
Version 1.0 was released in 2002. The project is now led by Moodle Pty Ltd
(Moodle HQ), together with a large community of developers and partners.

Moodle publishes two major releases per year. Some releases receive
long-term support (LTS): this is the case of Moodle 4.5, released in
October 2024 and supported until October 2027. The 5.x branch followed
from 2025 \cite{moodle_versions}. Both points matter for a final-year
project: the version used must include the AI subsystem (introduced in
4.5) and remain supported for the whole duration of the project.

\subsection{Technical architecture}

Moodle is a web application written in PHP. It relies on a relational
database (PostgreSQL, MySQL, MariaDB or Microsoft SQL Server) and
on a data directory (\texttt{moodledata}) that holds uploaded files, caches and
sessions. The whole platform is organized as \emph{plugins}: the core
provides shared services, and each visible feature is a plugin of a given
type. \Cref{tab-moodle-plugins} lists the plugin types most useful for our
project.

\begin{table}[htbp]
\centering
\caption{Main Moodle plugin types and their use in our project}
\label{tab-moodle-plugins}
\small
\begin{tabularx}{\textwidth}{@{}L{2.6cm}L{2.6cm}Y@{}}
\toprule
\textbf{Type} & \textbf{Prefix} & \textbf{Possible use in our project} \\
\midrule
Activity        & \texttt{mod}          & Activity dedicated to the tutor or to the adaptive path \\
Block           & \texttt{block}        & Recommendations shown on the course page \\
Local plugin    & \texttt{local}        & Event observers, additional web services \\
Question type   & \texttt{qtype}        & Questions generated and graded by an agent \\
Report          & \texttt{report}       & Dashboards for the teacher \\
AI provider     & \texttt{aiprovider}   & Connecting the AI subsystem to a chosen model \\
AI placement    & \texttt{aiplacement}  & New place in the interface where AI is offered \\
Log store       & \texttt{logstore}     & Exporting events to an LRS (xAPI) \\
Theme           & \texttt{theme}        & Look and feel of the platform \\
\bottomrule
\end{tabularx}
\end{table}

\subsection{Roles, permissions and contexts}

Moodle controls access through roles (manager, teacher, non-editing
teacher, student, guest) and capabilities, that is, elementary
permissions. A role is always assigned in a \emph{context}. Contexts form a
hierarchy: system, course category, course, activity module, block and
user \cite{moodle_contexts}. A teacher can therefore have rights in one
course and not in another.

This notion of context is organizational: it tells where the user is in
the structure of the platform. It differs from the learning context that
our framework must model (the learner's level, goals, history, time
constraints, device), which will be defined in \cref{chap-sota}. It is
still a usable source of information, and it imposes a constraint: the
framework must only access the data authorized in the context of the
request.

\subsection{Integration points}
\label{subsec-moodle-integration}

Moodle offers four mechanisms for exchanging data with an external service
without modifying its core. They are summarized in
\cref{fig-moodle-integration-points}.

\begin{itemize}
  \item \textbf{Web services.} Moodle exposes several hundred functions
        (reading courses, contents, grades and completions, creating
        questions, etc.). An external service calls them through REST with
        a token tied to an account and to a set of allowed functions
        \cite{moodle_ws}.
  \item \textbf{Events and logs.} Every significant action (viewing a
        resource, submitting an assignment, attempting a quiz) triggers an
        event. Events are recorded in the logs and can be caught by
        observers written in a plugin \cite{moodle_events}. They are the
        main source of the traces that will feed the learner model.
  \item \textbf{LTI.} Moodle can act as an LTI platform and launch an
        external tool from a course, or publish its own activities as a
        tool for other platforms. Version 1.3 and the LTI Advantage
        services are supported.
  \item \textbf{AI subsystem.} It connects language model providers and
        offers AI actions in the interface (\cref{subsec-moodle-ai}).
\end{itemize}

\begin{figure}[htbp]
\centering
\adjustbox{max width=\textwidth}{%
\begin{tikzpicture}[node distance=0.5cm and 1.6cm,
  b/.style={bloc,text width=3.4cm,minimum height=0.9cm}]
  \node[blocstrong,text width=3.4cm,minimum height=5.6cm] (moodle)
     {Moodle\\[4pt]\mdseries\footnotesize unmodified core\\ database\\ moodledata};
  \node[b,right=of moodle,yshift=2.1cm] (ws)  {REST web services};
  \node[b,below=of ws]                  (ev)  {Events and logs};
  \node[b,below=of ev]                  (lti) {LTI 1.3 and Advantage};
  \node[b,below=of lti]                 (ai)  {AI subsystem};
  \node[bloc,fill=white,draw=thaccent,text width=3.2cm,minimum height=5.6cm,
        right=1.6cm of ev.east,yshift=-0.7cm] (svc)
     {\textbf{External service}\\[4pt]\footnotesize
      reads courses, grades and traces; interface embedded in the
      course; calls language models};
  \foreach \n in {ws,ev,lti,ai}{
    \draw[arrow,<->] (moodle.east |- \n.west) -- (\n.west);
    \draw[arrow,<->] (\n.east) -- (svc.west |- \n.east);
  }
\end{tikzpicture}}
\caption{Moodle's integration points with an external service}
\label{fig-moodle-integration-points}
\end{figure}

\subsection{Adaptation tools available in Moodle}
\label{subsec-moodle-adaptation}

Moodle already offers several ways to adapt a course, all based on rules
written by hand by the teacher:

\begin{itemize}
  \item \textbf{Activity completion}: an activity is marked as complete
        according to criteria (viewed, minimum grade, submitted).
  \item \textbf{Restrict access}: a resource is only shown when conditions
        are met (date, grade obtained, completion of another activity,
        group membership).
  \item \textbf{Competencies and learning plans}: the teacher links
        competencies to activities and tracks how they are acquired.
  \item \textbf{Lesson and interactive quiz}: the Lesson activity branches
        according to the answers, and the quiz can run in interactive mode
        with several tries.
  \item \textbf{Analytics}: prediction models, such as identifying
        students at risk of dropping out, use the logs to alert the
        teacher \cite{moodle_analytics}.
\end{itemize}

These tools support branching adaptation, in the sense of adaptive
hypermedia \cite{brusilovsky2001}, but the teacher has to plan every path
in advance. The number of rules grows quickly with the number of
activities and learner profiles.

\subsection{The artificial intelligence subsystem}
\label{subsec-moodle-ai}

Since version 4.5, Moodle includes an AI subsystem in its core
\cite{moodle45,moodle_ai}. It is organized around three notions,
illustrated in \cref{fig-moodle-ai}:

\begin{itemize}
  \item \textbf{placements}: the places in the interface where users access
        AI, by default the text editor and course assistance;
  \item \textbf{actions}: what the AI can do, for example generate text or
        an image, summarize or explain content;
  \item \textbf{providers}: the language model services that carry out the
        actions, such as OpenAI or Azure AI, but also locally hosted models
        through Ollama \cite{moodle_ai}. A DeepSeek provider was added in
        Moodle 5.1 \cite{eledia2026}.
\end{itemize}

\begin{figure}[htbp]
\centering
\adjustbox{max width=\textwidth}{%
\begin{tikzpicture}[
  b/.style={bloc,text width=3.4cm,minimum height=0.8cm},
  ttl/.style={font=\footnotesize\bfseries,text=thaccent}]
  % three groups, centred on the same horizontal axis (y = 0)
  \node[b] (p1) at (0, 0.55) {Text editor};
  \node[b] (p2) at (0,-0.55) {Course assistance};
  \node[b] (a1) at (5.0, 1.1) {Generate text};
  \node[b] (a2) at (5.0, 0.0) {Generate an image};
  \node[b] (a3) at (5.0,-1.1) {Summarize, explain};
  \node[b] (f1) at (10.0, 1.1) {OpenAI};
  \node[b] (f2) at (10.0, 0.0) {Azure AI};
  \node[b] (f3) at (10.0,-1.1) {Ollama, others};
  \node[layer,fit=(p1)(p2)] (P) {};
  \node[layer,fit=(a1)(a3)] (A) {};
  \node[layer,fit=(f1)(f3)] (F) {};
  \foreach \g/\t in {P/Placements,A/Actions,F/Providers}
    \node[ttl,above=2pt] at (\g.north |- A.north) {\t};
  \draw[arrow] (P.east |- 0,0) -- (A.west |- 0,0);
  \draw[arrow] (A.east |- 0,0) -- (F.west |- 0,0);
  \node[below=0.5cm of A,font=\footnotesize,text=thgray,text width=11cm,align=center]
    {Configured by the administrator: activation, usage limits, system
     instructions; requests are logged};
\end{tikzpicture}}
\caption{Organization of Moodle's AI subsystem}
\label{fig-moodle-ai}
\end{figure}

Everything is disabled by default. The administrator chooses the allowed
providers, placements and actions, sets usage limits and can define a
system instruction added to every request. Usage is logged. This design
follows the AI principles published by Moodle, which put institutional
control and transparency first.

The subsystem is a solid foundation, but its role remains limited. It
handles isolated requests, with no memory of the learner's path. It does
not build a learner model, does not take the learner's context into
account, does not plan a learning path and does not coordinate several
specialized processes. Nor does it give teachers an analysis of their
materials based on the difficulties observed. These gaps match exactly the
scope of our framework.

% =====================================================================
\section{Limits of LMSs for adaptive learning}
\label{sec-limits}
% =====================================================================

LMSs were designed to organize teaching, not to support each learner.
\citet{coates2005} already observed that their use tends to reproduce
transmissive, uniform teaching. EDUCAUSE reports point the same way: users
mostly rely on the basic functions (sharing resources, assignments,
grades), and institutions want more open and more personalized
environments, which \citet{brown2015} call \emph{next generation digital
learning environments} \cite{dahlstrom2014,brown2015}. Our analysis
highlights six limits, summarized in \cref{tab-limits}.

\begin{enumerate}
  \item \textbf{A single path.} All enrolled learners usually see the same
        course in the same order. Adaptation exists, but it depends on
        hand-written rules (\cref{subsec-moodle-adaptation}).
  \item \textbf{No explicit learner model.} The LMS stores grades,
        completions and logs, but it does not build a representation of
        what the learner has mastered, of their recurring mistakes and of
        their goals.
  \item \textbf{Context reduced to the structure of the site.} The context
        known to the platform is limited to the course and the activity.
        The time of day, the workload, upcoming deadlines or the device in
        use are not exploited.
  \item \textbf{One-off AI assistance.} The built-in AI tools answer a
        specific request, with no coordination between functions and no
        link to the learner's tracking (\cref{subsec-moodle-ai}).
  \item \textbf{Little support for teachers to improve their materials.}
        Reports show who passed or failed, but they do not link the
        observed difficulties to the parts of the course that might
        explain them.
  \item \textbf{Scattered data.} Useful traces are spread across the logs,
        the gradebook, the question banks and sometimes an external LRS.
        Exploiting them requires integration work that the platform does
        not do.
\end{enumerate}

\begin{table}[htbp]
\centering
\caption{Limits of current LMSs and needs for adaptive learning}
\label{tab-limits}
\small
\begin{tabularx}{\textwidth}{@{}L{4.1cm}YY@{}}
\toprule
\textbf{Observed limit} & \textbf{Consequence} & \textbf{Need for our framework} \\
\midrule
Single path, hand-written rules & Struggling and advanced learners follow
  the same route & Automatic learning path planning \\
No learner model & Gaps are only detected after assessment & Continuous
  diagnosis of knowledge and mistakes \\
Context reduced to the site structure & Recommendations ignore the
  learner's actual situation & Explicit context model \\
One-off, isolated AI & No follow-up between requests, no coordination &
  Specialized agents coordinated by an orchestrator \\
Little help to improve materials & The same difficulties recur from one
  cohort to the next & Analysis of materials and proposals approved by the
  teacher \\
Scattered data & Costly exploitation of traces & A single integration
  layer (web services, events, xAPI) \\
\bottomrule
\end{tabularx}
\end{table}

\begin{summarybox}
LMSs, and Moodle in particular, provide a reliable infrastructure for
organizing teaching and an open integration layer (web services, events,
LTI 1.3, AI subsystem). However, they model neither the learner nor the
learner's context, and their AI tools remain one-off. An external
framework can fill these gaps by building on the platform's integration
points, without modifying its core.
\end{summarybox}

% =====================================================================
\section*{Conclusion}
\addcontentsline{toc}{section}{Conclusion}
% =====================================================================

This chapter presented online learning, the role and architecture of
LMSs, the standards that let them exchange data with other tools and the
main platforms on the market. The study of Moodle showed an open,
extensible and widely adopted platform, in particular in Algerian
universities, which has included a first AI subsystem since version 4.5.
It also showed that adaptation in Moodle still relies on hand-written
rules and that its built-in AI follows neither the learner nor the
learner's context.

Going further requires techniques able to model the learner, take their
context into account and coordinate several intelligent functions. The
next chapter studies these techniques: adaptive and context-aware
learning, large language models in education and multi-agent systems,
followed by the work that combines them.
